\documentclass[../tanulmany.tex]{subfiles}

\begin{document}

\subsection{The Debugging Paradox and the Glass Ceiling}

The rise of vibe coding marks a transformative era in software development, characterized by a fundamental shift from syntax-driven engineering to intention-driven orchestration. However, as this research has demonstrated, this transformation is accompanied by a critical phenomenon in the software development lifecycle.

The "Debugging Paradox" manifests in the widening gap between the ease of creation and the difficulty of maintenance. On one hand, LLMs have dramatically lowered the entry barrier to programming, allowing novices to build functional applications in hours that previously would have taken months to develop. \cite{finnie2022robots} On the other hand, this low barrier creates a deceptive "glass ceiling." Developers who rely exclusively on AI-generated code often build systems that exceed their own individual capacity to understand, debug, or refactor. 

As long as the AI's reasoning remains within the bounds of the developer's intent and the model's context window, the illusion of productivity holds. However, the moment the AI fails-whether through hallucination, logic errors in edge cases, or reaching the limits of its training data-the developer finds themselves in possession of a complex system they cannot maintain. This leads to a state of "technological helplessness" where the human component in the loop becomes the bottleneck, unable to provide the causal reasoning required to fix a failure in code they did not manually author.

\subsection{Long-Term Impact on Software Quality and Autonomy}

In the long term, the widespread adoption of vibe coding without the necessary pedagogical and architectural safeguards threatens to undermine both software quality and developer autonomy. If developers lose the ability to think from first principles, the industry may see a surge in "black box" systems-applications that work but whose internal logic is opaque even to their creators. 

Furthermore, the risk of skill decay suggests that a generation of programmers might become entirely dependent on proprietary, third-party LLMs, sacrificing their professional independence for short-term speed. \cite{macnamara2024does} In conclusion, while vibe coding is a powerful democratizing force, its success depends on our ability to maintain the human at the center of the logical loop. We must cultivate a culture of "critical orchestration," where the developer’s primary role is not just to prompt, but to audit, verify, and conceptually own every line of code generated by the machine. Only then can we ensure that the future of software development remains a discipline of rigorous engineering rather than one of mere statistical probability.

\end{document}